\section{Introduction}

AI for software engineering has made remarkable progress recently, becoming a notable success within generative AI.
Despite this, there are still many challenges that need to be addressed before automated software engineering reaches its full potential. With additional efforts, it should be possible to reach high levels of automation where humans can focus on the critical decisions of what to build and how to balance difficult tradeoffs while most routine development effort is automated away. Reaching this level of automation, however, will require substantial research and engineering efforts across academia and industry. \textbf{This paper provides an opinionated view of the tasks, challenges, and promising directions towards achieving this goal.}

Many existing surveys overlap with the topics that are discussed in this paper, such as AI for programming assistants \citep{liang2024large, sergeyuk2025using}, LLMs for software testing \citep{wang2024software}, using LLMs in low-resource and domain-specific languages \citep{joel2024survey}, automated program repair \citep{zhang2023survey} focus on automated program repair, and formal mathematical reasoning \citep{yang2024formal}. In addition, many papers discuss the current state, challenges, and future of AI for software engineering \citep{fan2023large, ozkaya2023application, wong2023natural, zheng2023survey, hou2024large, jin2024llms, wan2024deep, roychoudhury2025ai}. Our work draws inspiration from them, and we recommend that the reader consult with them for alternative perspectives. 

In this paper, our goal is threefold. In Sec. \ref{sec:main-tasks-milestones}, we provide a structured taxonomy of concrete tasks in AI for software engineering. In particular, we emphasize that there are many other tasks in SWE beyond code generation and code completion, encouraging research in these areas. Moving forward to Sec. \ref{sec:main-challenges}, we highlight four main challenges that today's models face, each cross-cutting and applicable to several tasks. In Sec. \ref{sec:main-paths}, we call for large open-source community efforts in the field and lay out a collection of promising research directions to address these challenges. The main text contains a summary of highlights, with a full version in the Appendix. We hope that after reading our paper, the reader can appreciate the progress AI for code has made, understand the shortcomings of today's state-of-the-art models, and come together to advance and shape the future of the field. 

% - AI coding has made progress last few years + is one of the success stories for genai

% - still a lot to go in terms of getting to fully automated swe

% - goal is to highlight the open problems and provide structure for the field going fwd

% - metrics, tasks, challenges framework

% - 2 purposes for paper existing: 1) highlight everything left to do, 2) provide this framework to structure the field

% alternative views: 
% - AI for code is a solved problem
% - we shouldn't be doing this because of economic reasons
% - AI safety: dangerous to have machines that can program themselves

